% ============================================================================
% Fichier  : partie4/P04_C08_conduire_audit.tex
% Titre    : Conduire un Audit de Protection des Données
% Auteur   : MINKA MI NGUIDJOÏ Thierry Emmanuel
% Version  : 3.0 -- novembre 2026
% Corrigés : corriges/Corr_C08.tex
% ============================================================================
% Sources : ISO 19011:2018 (principes et cycle, paraphrasés) ; ISO/IEC 27007
% (paraphrasée) ; auditologie (hal-05701007, hal-05701270) ; loi n° 2024/017,
% art. 35, 70 ; NRAC C06-C07 (quadrant, règle de décision).
% Faits nouveaux (bible) : plan d'audit de KORO AUDIT, six domaines,
% 2 juillet 2026.
% ============================================================================

\chapter{Conduire un Audit de Protection des Données}
\label{chap:audit}
\chaptermark{Conduire un audit}

\epigraph{%
    « Un audit ne dit pas si l'organisation est bonne.\\
    Il dit ce qui est établi, et sur quelle preuve. »%
}{Principe du cours}

\niveauChapitre{Fondamental}{%
    Chapitres~\ref{chap:donnees} et~\ref{chap:principes}~; méthode générale
    de l'audit (UE MSN-2S3-SEC08)
}

\begin{boxobjectifs}
\begin{itemize}
    \item Distinguer audit de conformité légale, audit contre une norme et
          audit de préparation.
    \item Rédiger le mandat, les critères et le périmètre d'une mission.
    \item Choisir des preuves selon leur force et documenter chaque constat
          dans un papier de travail.
    \item Décomposer la conformité en domaines et en exigences cotées
          \OUI{}, \NON{} ou \NEK{}.
    \item Lire les mesures $\mu$, $U$ et $H$ et appliquer une règle de
          décision écrite à l'avance.
    \item Protéger les données personnelles que l'audit lui-même manipule.
\end{itemize}
\end{boxobjectifs}

\bigskip

% ============================================================================
\section*{Avant de commencer : « montrez-nous la sécurité »}
% ============================================================================

Yaoundé, jeudi 2~juillet~2026. Dans les bureaux de KORO AUDIT \&
CONFORMITÉ, Mme~BELLO prépare la mission de NYANGA CRÉDIT avec son associé
juriste, M.~NGANDO. Le DSI de l'établissement a envoyé la veille un dossier
de quarante pages~: schémas du réseau, règles du pare-feu, rapport du test
d'intrusion de 2025. « Voilà notre conformité », écrit-il.

M.~NGANDO le parcourt et le referme. Le dossier ne dit rien des
consentements, des durées, des droits des personnes, des contrats des
prestataires ni des transferts. Il répond à une question que personne n'a
posée. Mme~BELLO commence par écrire ce que la mission doit établir, avec
quels critères, sur quel périmètre, et ce qu'elle ne couvrira pas.

\begin{boxalert}
\textbf{Violation typique.} Confondre l'audit de protection des données avec
un audit de sécurité informatique, et conclure à la conformité parce que les
systèmes sont bien protégés.

\smallskip
\textbf{Règle~:} les critères de l'audit sont les obligations de la loi et,
le cas échéant, d'une norme désignée. La sécurité n'en est qu'un domaine
parmi six~; des données bien gardées peuvent avoir été collectées sans
droit.
\end{boxalert}

% ============================================================================
\section{Trois Sortes d'Audit}
\label{sec:c03-types}
% ============================================================================

\begin{tableaucours}{L{3.2cm}L{3.8cm}L{5.0cm}}{%
    \tableauHeader{Audit} &
    \tableauHeader{Critères} &
    \tableauHeader{Usage}
}
De conformité légale & Loi n°~2024/017, textes d'application, autres lois
applicables & Dire si l'organisation respecte ses obligations~; préparer un
contrôle ou un contentieux \\
Contre une norme & ISO/IEC~27701, ISO/IEC~27001 & Préparer ou maintenir une
certification (Chapitre~\ref{chap:normes}) \\
De préparation & Loi et bonnes pratiques & Mesurer l'écart avant une
échéance~; prioriser les chantiers \\
\end{tableaucours}
\vspace{-0.8em}
\begin{center}{\small\itshape Tableau~8.1. Trois audits, trois jeux de critères}\end{center}

Ces audits peuvent être conduits par l'organisation elle-même (audit
interne), par un cabinet qu'elle mandate (NYANGA, MBOLO), ou par un tiers qui
certifie. Le contrôle de l'Autorité de protection, prévu par la loi
(article~53) et dont les modalités relèvent de textes réglementaires
(article~35), n'est pas un audit au sens de ce chapitre~: c'est un contrôle
administratif, qui peut conduire à des sanctions (Chapitre~\ref{chap:loi}).

% ============================================================================
\section{Les Principes de l'Auditeur}
\label{sec:c03-principes}
% ============================================================================

La norme ISO~19011~\cite{iso19011_2018} pose sept principes pour tout audit
de système de management. Ils s'appliquent sans changement à l'audit de
protection des données~; on les résume ici en paraphrase.

\begin{enumerate}
    \item \textbf{Intégrité}~: honnêteté, diligence, respect des lois.
    \item \textbf{Présentation impartiale}~: rendre compte avec exactitude,
          y compris des difficultés rencontrées.
    \item \textbf{Conscience professionnelle}~: le soin et le jugement
          proportionnés à l'importance de la mission.
    \item \textbf{Confidentialité}~: la protection des informations
          obtenues, et ici des données personnelles consultées.
    \item \textbf{Indépendance}~: l'absence de conflit d'intérêts, fondement
          de l'objectivité des conclusions.
    \item \textbf{Approche fondée sur la preuve}~: des conclusions
          vérifiables, tirées d'un échantillon défini.
    \item \textbf{Approche par les risques}~: concentrer l'effort là où
          l'écart serait le plus grave pour les personnes.
\end{enumerate}

Le quatrième principe prend ici un relief particulier. Pour vérifier un
consentement, l'auditeur ouvre un dossier client~; pour tester les droits,
il lit des réponses adressées à des personnes~; pour contrôler les
journaux, il voit qui a consulté quoi. L'audit est lui-même un traitement~:
il a une finalité (la mission), une durée (celle de la conservation des
papiers de travail) et des destinataires (le commanditaire). L'auditeur
consulte sur place, copie le moins possible, masque les identités dans ses
papiers et détruit les copies à l'échéance prévue au mandat
(section~\ref{sec:c08-deontologie}).

% ============================================================================
\section{La Déontologie de l'Auditeur}
\label{sec:c08-deontologie}
% ============================================================================

L'auditeur de la protection des données est soumis aux principes communs de
l'audit de la section précédente. Son objet ajoute trois exigences.

\subsection{Minimiser les données de l'audit}

L'auditeur ne collecte que les données nécessaires à la preuve. Il consulte
sur place plutôt que de copier, relève des indicateurs plutôt que des
contenus, masque les identités dans ses papiers de travail (« client~A,
dossier 2018-0412 »), conserve les copies indispensables dans un espace
chiffré et les détruit à la date prévue au mandat. Un auditeur qui
emporterait un extrait du fichier clients sur une clé non chiffrée
commettrait l'écart même qu'il est chargé de relever.

\subsection{L'indépendance}

L'indépendance interdit d'auditer son propre travail. KORO AUDIT a proposé
au CENTRE MÉDICAL MBOLO, en juin, de rédiger ses procédures de protection
des données~; s'il accepte, il ne pourra pas, en octobre, auditer ces
mêmes procédures sans menacer son objectivité. Mme~BELLO a refusé la
mission de rédaction et conservé la mission d'audit.

\subsection{Ce que l'auditeur découvre}

L'auditeur peut découvrir des faits graves~: un agent qui vend des fichiers,
un traitement manifestement délictueux. Il en rend compte au commanditaire,
par écrit et sans délai, comme son mandat doit le prévoir. Il ne mène pas
d'enquête qui excéderait sa mission, et il ne saisit pas lui-même la
justice, sauf obligation légale qui lui serait propre. Le
Chapitre~\ref{chap:violations} montre comment un tel fait devient une
violation à traiter, et le cours de forensique comment il devient une
enquête.

% ============================================================================
\section{Mandat, Critères, Périmètre}
\label{sec:c03-mandat}
% ============================================================================

\begin{boxprocedure}
\textbf{Préparer une mission d'audit de protection des données.}
\begin{enumerate}
    \item Obtenir un mandat écrit~: commanditaire, objectif, accès accordés,
          destinataires du rapport, date de remise.
    \item Fixer les critères~: articles de la loi retenus, textes
          d'application publiés, norme éventuelle, politiques internes.
    \item Délimiter le périmètre~: entités, sites, traitements, période~;
          écrire ce qui est exclu, et pourquoi.
    \item Décomposer en domaines et en exigences, chacune formulée en
          question fermée.
    \item Choisir, pour chaque exigence, le test et la preuve attendue.
    \item Prévoir le traitement des données de l'audit~: consultation sur
          place, masquage, durée de conservation des papiers.
    \item Fixer à l'avance la règle de décision qui permettra de conclure à
          partir des cotes.
\end{enumerate}
\end{boxprocedure}

Les critères doivent être \textbf{publiés et applicables}. Un critère qui
renvoie à un texte réglementaire non publié, comme la durée de conservation
fixée « suivant les modalités fixées par voie réglementaire »
(article~13), ne peut pas être appliqué tel quel~: l'auditeur le remplace
par un critère de substitution déclaré (« une durée fixée, justifiée et
appliquée par l'organisation »), ou il cote \NEK{} en précisant que
l'incertitude tient au critère et non à la preuve. Les deux sources
d'incertitude se distinguent dans le rapport~: un \NEK{} \textit{de
preuve} se lève en cherchant une pièce, un \NEK{} \textit{de critère} ne se
lève qu'avec un texte.

% ============================================================================
\section{Les Preuves et leur Force}
\label{sec:c03-preuves}
% ============================================================================

\begin{tableaucours}{cL{3.4cm}L{4.4cm}L{3.4cm}}{%
    \tableauHeader{Rang} &
    \tableauHeader{Type de preuve} &
    \tableauHeader{Exemple} &
    \tableauHeader{Limite}
}
1 & Réexécution, test & Demande d'accès test suivie jusqu'à la réponse &
Coûteuse, à encadrer \\
2 & Trace produite par le système & Journal d'accès, horodatage de
consentement & Dépend de la configuration \\
3 & Document signé et daté & Contrat de sous-traitance, autorisation &
Prouve l'engagement, pas son exécution \\
4 & Observation & Affichage en agence, armoire ouverte & Instantanée \\
5 & Déclaration & Entretien, questionnaire & À corroborer \\
\end{tableaucours}
\vspace{-0.8em}
\begin{center}{\small\itshape Tableau~8.2. Hiérarchie indicative des preuves}\end{center}

Une déclaration seule ne permet jamais de coter \OUI{}~; elle oriente la
recherche d'une preuve de rang supérieur. Un document de rang~3 établit
l'existence d'une règle, pas son application~: un contrat qui impose des
mesures de sécurité au prestataire ne prouve pas que le prestataire les
applique.

Chaque constat se documente dans un \termecle{papier de travail}~:
l'exigence, le critère, le test, l'échantillon, la preuve (référence de la
pièce), la cote, la date, l'auteur et le réviseur. Le papier de travail
permet à un autre auditeur de refaire le raisonnement, et à l'audité de
contester un constat précis plutôt qu'une impression.

\begin{boxinfo}
\textbf{La demande test.} Pour vérifier le traitement des demandes de
droits, l'auditeur peut adresser, avec l'accord écrit du commanditaire, une
demande d'accès au nom d'un membre de l'équipe qui est réellement client.
Une demande au nom d'une personne fictive, ou d'un client à son insu,
serait une collecte déloyale. Le mandat doit prévoir ce test.
\end{boxinfo}

% ============================================================================
\section{Décomposer et Coter}
\label{sec:c03-decomposer}
% ============================================================================

La conformité à une loi de soixante-quinze articles ne se cote pas d'un
bloc. On la décompose en \termecle{domaines}, puis en \termecle{exigences},
chacune posée comme une question fermée à laquelle une preuve peut répondre
\OUI{}, \NON{} ou \NEK{}.

\begin{tableaucours}{cL{3.6cm}L{3.4cm}L{4.4cm}}{%
    \tableauHeader{Code} &
    \tableauHeader{Domaine} &
    \tableauHeader{Articles} &
    \tableauHeader{Exemple d'exigence}
}
D1 & Recensement et fondements & 9--14, 19, 29 & Le registre existe-t-il~? \\
D2 & Droits des personnes & 21, 37--47 & Les demandes reçoivent-elles une
réponse~? \\
D3 & Sécurité et confidentialité & 15, 24, 25, 27 & Les accès sont-ils
tracés~? \\
D4 & Sous-traitants et transferts & 16, 26, 30--32 & Chaque prestataire
a-t-il un contrat conforme~? \\
D5 & Violations & 22 & Les violations sont-elles notifiées~? \\
D6 & Gouvernance & 27(2), 33, 35 & L'évaluation d'impact est-elle faite
quand le risque est élevé~? \\
\end{tableaucours}
\vspace{-0.8em}
\begin{center}{\small\itshape Tableau~8.3. Les six domaines d'un audit de protection des données}\end{center}

Une exigence formulée avec « et » ou « ou » se scinde~: « le registre existe
et il est à jour » contient deux questions, qui peuvent recevoir deux cotes
différentes. La règle d'atomicité évite qu'un \NON{} ne désigne deux défauts
à la fois.

% ============================================================================
\section{Lire les Mesures, Conclure}
\label{sec:c03-conclure}
% ============================================================================

\begin{boxdef}
Pour un ensemble de $n$ exigences cotées, avec $n_O$ réponses \OUI{}, $n_N$
\NON{} et $n_K$ \NEK{}~:

\smallskip
\termecle{Taux de conformité}~: $\mu = n_O / n$.
\termecle{Taux d'incertitude}~: $U = n_K / n$.
\termecle{Entropie}~: $H = -\sum p \log_2 p$, sur les trois cotes, en bits,
entre~0 et $\log_2 3 \approx 1{,}585$.
\end{boxdef}

\begin{boxprocedure}
\textbf{La règle de décision du cours.}
\begin{enumerate}
    \item \textbf{$U$ d'abord.} Si $U \geq 0{,}20$, le domaine est opaque~:
          aucune conclusion de conformité. On recherche les preuves, ou l'on
          conclut sur l'opacité elle-même.
    \item \textbf{$\mu$ ensuite.} Si $U < 0{,}20$, le domaine est clair~:
          conforme si $\mu \geq 0{,}5$ et sans \NON{} sur une exigence
          pénalement sanctionnée~; non conforme sinon.
    \item \textbf{$H$ pour contrôler.} Une entropie élevée avec $U$ faible
          signale un désaccord entre \OUI{} et \NON{}~: on décompose avant
          de conclure.
    \item \textbf{Les cotes prévalent.} Un seul \NON{} grave (collecte
          déloyale, transfert illicite, données sensibles sans
          consentement) suffit à motiver une recommandation prioritaire,
          quelles que soient les mesures.
\end{enumerate}
\end{boxprocedure}

Les seuils sont ceux du cours~; chaque mission peut en fixer d'autres, à
condition de les écrire dans le plan d'audit avant de coter. La quatrième
règle est propre à la protection des données~: la loi sanctionne
pénalement certains écarts (articles~63 à~71), et un taux de conformité
élevé ne compense pas une infraction.

% ============================================================================
\section{Le Plan d'Audit de NYANGA}
\label{sec:c03-nyanga}
% ============================================================================

\begin{boxscenario}{Mission NYANGA. Le plan d'audit de KORO AUDIT, 2 juillet 2026}
\textbf{Mandat.} Conseil d'administration de NYANGA CRÉDIT, 30~juin~2026~:
« dire si l'établissement respecte la loi n°~2024/017 et recommander les
mesures nécessaires ». Rapport au conseil le 25~septembre.

\textbf{Critères.} Les articles~2 à~52 de la loi n°~2024/017~; les
articles~26, 41 à~48 et~74 de la loi n°~2010/012~; les textes d'application
publiés à la date de la mission (aucun identifié)~; la politique de sécurité
interne de 2025. La carte d'applicabilité du Chapitre~\ref{chap:loi2010} est
jointe au plan.

\textbf{Périmètre.} Siège et trois agences tirées au sort (Dschang,
Mbouda, Foumban)~; les onze familles de données du
Chapitre~\ref{chap:donnees}~; période du 1\textsuperscript{er}~janvier~2025
au 31~août~2026. \textbf{Exclusions déclarées~:} l'enquête pénale sur
l'incident de mars~2025 et le contrôle de l'éditeur du logiciel bancaire
dans ses propres locaux, à Abidjan.

\textbf{Données de l'audit.} Consultation dans les locaux de NYANGA~;
papiers de travail avec identités masquées~; conservation trois ans, puis
destruction attestée.

\textbf{Règle de décision.} Celle du cours, seuils inchangés.
\end{boxscenario}

\begin{boxlecture}{Le domaine D1 après la première semaine}
\begin{tableaucours}{L{7.8cm}cL{3.2cm}}{%
    \tableauHeader{Exigence} &
    \tableauHeader{Cote} &
    \tableauHeader{Preuve}
}
Les familles de données sont recensées & \OUI{} & Inventaire du 7~juillet \\
Un registre des activités de traitement est tenu & \NON{} & Aucun registre \\
Chaque traitement a un fondement valable & \NON{} & Formulaires, cahiers \\
L'autorisation de l'Autorité est demandée & \NON{} & Attestation \\
Une durée de conservation est fixée & \NON{} & Aucune règle \\
L'information de l'article~21 est fournie & \NON{} & Formulaire 2016 \\
L'identification des clients repose sur une obligation légale & \NEK{} &
Texte non produit \\
La vidéosurveillance repose sur un fondement valable & \NEK{} & Fondement
incertain \\
\end{tableaucours}

\medskip
\textbf{Décompte.} Un \OUI{}, cinq \NON{}, deux \NEK{} sur huit.
\hfill\lecture{0{,}125}{0{,}25}{1{,}30}

\medskip
\textbf{Lecture.} $U = 0{,}25 \geq 0{,}20$~: selon la première règle, le
domaine est opaque et ne permet aucune conclusion de conformité. Mais la
quatrième règle s'applique déjà~: cinq écarts sont établis, dont
l'absence d'autorisation, sanctionnée par l'article~55. La conclusion
provisoire est une non-conformité établie sur cinq exigences, avec deux
exigences à instruire.
\end{boxlecture}

% ============================================================================
\section{Ce que la Méthode demande à chaque Métier}
\label{sec:c03-metiers}
% ============================================================================

\begin{boxmetier}{Auditeur de la protection des données}
Écrire le mandat, les critères, le périmètre et la règle de décision avant
le premier entretien. Distinguer dans le rapport les \NEK{} de preuve et les
\NEK{} de critère.
\end{boxmetier}

\begin{boxmetier}{Responsable de traitement}
Donner à l'auditeur un accès réel aux pièces, et nommer un interlocuteur
par domaine. L'entrave à un contrôle de l'Autorité est un délit
(article~70)~; l'entrave à un audit interne prive surtout l'organisation
de sa propre information.
\end{boxmetier}

\begin{boxmetier}{RSSI}
Préparer les preuves de sécurité sous une forme exploitable~: journaux
exportés et datés, listes d'accès, rapports de sauvegarde. Un schéma de
réseau ne prouve aucune mesure de protection des données.
\end{boxmetier}

\begin{boxresume}
\begin{itemize}
    \item L'audit de protection des données a pour critères la loi et, le
          cas échéant, une norme~; la sécurité n'en est qu'un domaine.
    \item Les sept principes de l'ISO~19011 s'appliquent~; la
          confidentialité impose de protéger les données que l'audit
          manipule.
    \item Un critère renvoyant à un texte non publié produit un \NEK{} de
          critère, distinct d'un \NEK{} de preuve.
    \item La conformité se décompose en six domaines et en exigences
          atomiques cotées \OUI{}, \NON{} ou \NEK{}.
    \item La règle de décision lit $U$ d'abord, $\mu$ ensuite, et un seul
          écart grave suffit à motiver une recommandation prioritaire.
\end{itemize}
\end{boxresume}

\begin{boxinfo}
\textbf{Points de contrôle.} Vous avez compris ce chapitre si vous savez~:
\begin{itemize}
    \item rédiger le mandat et le périmètre d'une mission~;
    \item classer des preuves selon leur force~;
    \item appliquer la règle de décision à un domaine coté.
\end{itemize}
\end{boxinfo}

\begin{boxexo}[title={\ding{45}\quad Exercice~8.1. Quelle preuve \niveaubase}]
Pour chaque exigence, proposez la preuve la plus forte possible~: (a) les
salariés ont été informés du traitement de leurs données~; (b) le
prestataire de recouvrement efface les dossiers soldés~; (c) les accès aux
dossiers clients sont limités aux agents de l'agence~; (d) les demandes
d'accès reçoivent une réponse.
\end{boxexo}

\begin{boxexo}[title={\ding{45}\quad Exercice~8.2. Rendre atomique \niveaumoyen}]
Réécrivez en exigences atomiques~: (a) « le registre existe et il est tenu à
jour par un responsable désigné »~; (b) « les données sont chiffrées ou
pseudonymisées »~; (c) « les contrats des prestataires prévoient la
sécurité, la confidentialité et la restitution des données ».
\end{boxexo}

\begin{boxexo}[title={\ding{45}\quad Exercice~8.3. Conclure sur D2 \niveaumoyen}]
Le domaine D2 de NYANGA compte dix exigences~: 2~\OUI{}, 5~\NON{}, 3~\NEK{}.
Calculez $\mu$, $U$ et $H$, appliquez la règle de décision et rédigez la
phrase de conclusion du domaine.
\end{boxexo}

\begin{boxexo}[title={\ding{45}\quad Exercice~8.4. L'audit comme traitement \niveauexpert}]
Rédigez la clause du mandat qui encadre les données personnelles consultées
par l'auditeur~: finalité, accès, copies, masquage, conservation,
destruction. Dites comment l'auditeur prouvera lui-même qu'il l'a respectée.
\end{boxexo}

\begin{boxexo}[title={\ding{45}\quad Exercice~8.5. Une découverte \niveauexpert}]
Pendant l'audit de MBOLO, un auditeur junior découvre qu'un agent exporte
chaque semaine la liste des patientes de la maternité vers une adresse de
messagerie personnelle. Rédigez la note au commanditaire, et dites ce que
l'auditeur doit faire et ne pas faire.
\end{boxexo}

\bigskip

\begin{boxinfo}
\textbf{Pour aller plus loin.}
\begin{itemize}
    \item L'ISO~19011~\cite{iso19011_2018} et l'ISO/IEC~27007~\cite{iso27007},
          pour la conduite des audits.
    \item T.~E.~Minka Mi Nguidjoï, « Vers l'auditologie »~\cite{minka2026auditologie},
          pour la décomposition et les mesures d'incertitude.
\end{itemize}
\end{boxinfo}

\begin{boxinfo}
\textbf{Transition.} La première exigence de D1 est le registre. Comment
recenser tous les traitements d'une organisation, y compris ceux que la
direction ignore~? C'est l'objet du Chapitre~\ref{chap:cartographie}.
\end{boxinfo}
